% \chapter*{Resumen}
% \addcontentsline{toc}{chapter}{Resumen ejecutivo}
% \justifying
% El desempeño de un cable eléctrico subterráneo está limitado por la temperatura alcanzada por su conductor, la cual depende de las pérdidas internas y de la capacidad del entorno para evacuar calor. Los métodos normativos de ampacidad son indispensables como referencia de ingeniería, pero representan el terreno mediante resistencias térmicas equivalentes y supuestos simplificados; en campo, la humedad, el secado, los estratos, el lecho de asiento y el suelo de relleno producen una conductividad térmica espacialmente variable. Esta diferencia puede ocultar puntos calientes o inducir estimaciones conservadoras de capacidad \parencites[4--6]{enescu2021}[41--42]{cigre2025}.

% La investigación propone diseñar y evaluar, bajo el paradigma de ciencia del diseño (DSR), un artefacto computacional que incorpore la ecuación de conducción de calor, las condiciones de contorno e interfaz y la conductividad térmica espacialmente variable $\kx$ dentro del entrenamiento de una red neuronal informada por física (PINN). El artefacto será demostrado en casos analíticos, benchmarks FEM y escenarios documentados de cables enterrados; su evaluación considerará exactitud térmica, conservación de energía, robustez y capacidad para estimar la ampacidad. La secuencia metodológica adapta las actividades de identificación del problema, objetivos de solución, diseño, demostración, evaluación y comunicación de DSRM, con ciclos formativos de refinamiento durante el desarrollo \parencites[55--56]{peffers2007}[200--202]{delport2024}.

% Se espera producir un modelo lógico, un prototipo reproducible, un conjunto de escenarios y principios de diseño transferibles a la clase de problemas de conducción térmica 2D multimaterial. La contribución no se reducirá al software: incluirá criterios de representación de heterogeneidad, procedimientos de verificación y evidencia sobre cuándo la aproximación homogénea puede sesgar la temperatura máxima o la corriente admisible \parencites[343--350]{gregor2013}[93--99]{cigre2025}.

% \noindent\textbf{Palabras clave:} cable subterráneo; ampacidad; heterogeneidad térmica; conductividad del suelo; red neuronal informada por física; ciencia del diseño.
